
% =========================================================
\section{The flip-to-switching separation law}
\label{sec:separation-law}
% =========================================================

We now determine the parameter value at which the smooth
period-two branch created at the flip bifurcation first reaches the
gradient-clipping boundary.

Recall that
\[
\mu=\eta-\eta_f,
\qquad
\eta_f=\frac{2}{\lambda_*},
\]
and let
\[
\mathcal O_2(\mu)
=
\{x_-(\mu),x_+(\mu)\}
\]
denote the smooth period-two branch constructed in
Theorem~\ref{thm:period-two-branch}.

By Proposition~\ref{prop:exact-gradient-balance}, the two points of
the cycle have exactly opposite gradients:
\[
\nabla V(x_+(\mu))
=
-\nabla V(x_-(\mu)).
\]
Hence they have a common clipping exposure
\begin{equation}
\mathscr E(\mu)
:=
\|\nabla V(x_+(\mu))\|_\diamond
=
\|\nabla V(x_-(\mu))\|_\diamond.
\label{eq:section5-exposure}
\end{equation}

The first switching contact therefore occurs when the single scalar
condition
\begin{equation}
\mathscr E(\mu)=c
\label{eq:contact-equation}
\end{equation}
is satisfied.

The main purpose of this section is to solve
\eqref{eq:contact-equation} asymptotically as $c\to0^+$.

Quadratic tangency between period-doubling and border-collision
loci is known for generic piecewise-smooth maps
\cite{SimpsonMeiss2009}.
Accordingly, the exponent two obtained below is not claimed as a
generic new border-collision phenomenon.

The contribution here is the explicit gradient-induced coefficient,
the simultaneous nature of the two-point contact, and the improved
even-power remainder resulting from the exact gradient-balance
structure.

% ---------------------------------------------------------
\subsection{Factorization of the gradient exposure}
\label{subsec:exposure-factorization}
% ---------------------------------------------------------

The square-root law in
Theorem~\ref{thm:exposure-expansion} can be strengthened into a
factorized form that is particularly convenient for inversion.

\begin{lemma}[Exposure factorization]
\label{lem:exposure-factorization}
Assume {\rm (H1)--(H5)}. There exist $\mu_0>0$ and a positive
$C^1$ function
\[
\beta:[0,\mu_0)\to(0,\infty)
\]
such that
\[
E(\mu)=\sqrt{\mu}\,\beta(\mu),
\qquad
0\leq\mu<\mu_0,
\]
with
\[
\beta(0)
=
B_*
=
\lambda_*A_*\|v_*\|_\diamond.
\]
Consequently,
\[
B_*^2
=
\frac{3\lambda_*^4\|v_*\|_\diamond^2}{\Gamma_*}.
\]
\end{lemma}

\begin{proof}
Introduce the signed square-root parameter
\[
\varepsilon^2=\mu.
\]
By Proposition~\ref{prop:refined-cycle-expansion}, the
half-difference of the period-two orbit satisfies
\[
d(-\varepsilon)=-d(\varepsilon)
\]
and
\[
d(\varepsilon)
=
A_*\varepsilon v_*
+
O(\varepsilon^3).
\]

The signed parametrization of
Section~\ref{subsec:square-root-parity} shows that $d$ is a $C^3$
odd function of $\varepsilon$.  Hence it can be written locally as
\begin{equation}
d(\varepsilon)
=
\varepsilon q(\varepsilon^2),
\label{eq:d-factorization}
\end{equation}
where $q$ is a $C^1$ function and
\begin{equation}
q(0)=A_*v_*.
\label{eq:q-zero}
\end{equation}

By the exact gradient-balance identity
\eqref{eq:exact-exposure-d},
\[
\mathscr E(\mu)
=
\frac{2}{\eta_f+\mu}
\|d(\sqrt{\mu})\|_\diamond.
\]
Substituting \eqref{eq:d-factorization} gives
\[
\mathscr E(\mu)
=
\sqrt{\mu}
\frac{2}{\eta_f+\mu}
\|q(\mu)\|_\diamond.
\]
Thus define
\begin{equation}
\beta(\mu)
:=
\frac{2}{\eta_f+\mu}
\|q(\mu)\|_\diamond.
\label{eq:beta-definition}
\end{equation}

Since
\[
q(0)=A_*v_*\neq0
\]
and the norm $\|\cdot\|_\diamond$ is smooth away from the
origin by {\rm (H5)}, the function $\beta$ is $C^1$ after
reducing $\mu_0$ if necessary.

Furthermore,
\begin{align}
\beta(0)
&=
\frac{2}{\eta_f}
\|A_*v_*\|_\diamond
\nonumber\\
&=
\lambda_*A_*\|v_*\|_\diamond
=
B_*,
\end{align}
because
\[
\frac{2}{\eta_f}=\lambda_*.
\]

Finally,
\[
A_*^2
=
\frac{3\lambda_*^2}{\Gamma_*}
\]
by Proposition~\ref{prop:amplitude-identity}. Hence
\[
B_*^2
=
\lambda_*^2A_*^2\|v_*\|_\diamond^2
=
\frac{
3\lambda_*^4\|v_*\|_\diamond^2
}{
\Gamma_*
}.
\]
\end{proof}

\begin{remark}
\label{rem:exposure-factorization-importance}
The factorization
\[
\mathscr E(\mu)
=
\sqrt{\mu}\,\beta(\mu)
\]
is stronger than a formal leading-order expansion.
It separates the universal square-root bifurcation scale from a
regular geometric factor $\beta(\mu)$.

This form will allow the first-contact problem to be reduced to a
regular implicit-function problem in the variable $c^2$.
\end{remark}

% ---------------------------------------------------------
\subsection{Monotonicity and uniqueness of the first contact}
\label{subsec:contact-uniqueness}
% ---------------------------------------------------------

We next show that the exposure increases strictly along the newborn
period-two branch.

\begin{proposition}[Local monotonicity of the exposure]
\label{prop:exposure-monotonicity}
There exists $\mu_1>0$ such that
\begin{equation}
\frac{d}{d\mu}\mathscr E(\mu)>0
\qquad
\text{for all }
0<\mu<\mu_1.
\label{eq:exposure-monotonicity}
\end{equation}
Consequently, $\mathscr E$ is strictly increasing on
$(0,\mu_1)$.
\end{proposition}

\begin{proof}
By Lemma~\ref{lem:exposure-factorization},
\[
\mathscr E(\mu)
=
\sqrt{\mu}\,\beta(\mu),
\]
where
\[
\beta(0)=B_*>0.
\]
For $\mu>0$,
\begin{equation}
\mathscr E'(\mu)
=
\frac{\beta(\mu)}{2\sqrt{\mu}}
+
\sqrt{\mu}\,\beta'(\mu).
\label{eq:exposure-derivative}
\end{equation}

Since $\beta$ is continuous and $\beta(0)=B_*>0$, there exists
$\mu_1>0$ such that
\[
\beta(\mu)\geq\frac{B_*}{2}
\]
for $0\leq\mu<\mu_1$.
Moreover, $\beta'$ is bounded on this interval. Therefore the first
term in \eqref{eq:exposure-derivative}, which is of order
$\mu^{-1/2}$, dominates the second term for sufficiently small
$\mu>0$.

Hence
\[
\mathscr E'(\mu)>0
\]
after reducing $\mu_1$ if necessary.
\end{proof}

We may therefore define the first-contact parameter unambiguously.

\begin{definition}[First switching-contact parameter]
\label{def:first-contact}
For sufficiently small $c>0$, let
$\mu_{\mathrm{sc}}(c)$ denote the unique positive solution of
\begin{equation}
\mathscr E(\mu_{\mathrm{sc}}(c))=c.
\label{eq:mu-sc-definition}
\end{equation}
The corresponding step size is
\begin{equation}
\eta_{\mathrm{sc}}(c)
:=
\eta_f+\mu_{\mathrm{sc}}(c).
\label{eq:eta-sc-definition}
\end{equation}
\end{definition}

% ---------------------------------------------------------
\subsection{The quadratic separation law}
\label{subsec:quadratic-separation}
% ---------------------------------------------------------

We now state the central result of the paper.

\begin{theorem}[Flip-to-switching separation law]
\label{thm:quadratic-separation}
Assume \textbf{(H1)--(H5)}.
There exists $c_0>0$ such that for every
\[
0<c<c_0
\]
the first-contact parameter
$\eta_{\mathrm{sc}}(c)$ is well defined and unique.

As $c\to0^+$,
\begin{equation}
\boxed{
\eta_{\mathrm{sc}}(c)-\eta_f
=
K_\diamond c^2
+
O(c^4)
}
\label{eq:main-separation-law}
\end{equation}
where
\begin{equation}
\boxed{
K_\diamond
=
\frac{1}{
\lambda_*^2A_*^2
\|v_*\|_\diamond^2
}
}
\label{eq:K-A}
\end{equation}
or, equivalently,
\begin{equation}
\boxed{
K_\diamond
=
\frac{\ell_*}{
\lambda_*^3
\|v_*\|_\diamond^2
}
}
\label{eq:K-ell}
\end{equation}
and
\begin{equation}
\boxed{
K_\diamond
=
\frac{\Gamma_*}{
3\lambda_*^4
\|v_*\|_\diamond^2
}.
}
\label{eq:K-Gamma}
\end{equation}

In particular,
\begin{equation}
\lim_{c\to0^+}
\frac{
\eta_{\mathrm{sc}}(c)-\eta_f
}{
c^2
}
=
K_\diamond>0.
\label{eq:K-limit}
\end{equation}
\end{theorem}

\begin{proof}
By Proposition~\ref{prop:exposure-monotonicity}, the exposure
$E(\mu)$ is strictly increasing for all sufficiently small
$\mu>0$. Since
\[
E(\mu)\to0
\qquad\text{as}\qquad
\mu\to0^+,
\]
continuity and strict monotonicity imply that, for every
sufficiently small $c>0$, there exists a unique
$\mu_{\mathrm{sc}}(c)>0$ satisfying
\[
E(\mu_{\mathrm{sc}}(c))=c.
\]

By Theorem~\ref{thm:exposure-expansion},
\[
E(\mu)
=
B_*\sqrt{\mu}
+
O(\mu^{3/2})
=
\sqrt{\mu}\bigl(B_*+O(\mu)\bigr).
\]
Therefore, at $\mu=\mu_{\mathrm{sc}}(c)$,
\[
c^2
=
\mu_{\mathrm{sc}}(c)
\bigl(B_*+O(\mu_{\mathrm{sc}}(c))\bigr)^2,
\]
and hence
\[
c^2
=
B_*^2\mu_{\mathrm{sc}}(c)
+
O\bigl(\mu_{\mathrm{sc}}(c)^2\bigr).
\]

Since
\[
\frac{E(\mu)}{\sqrt{\mu}}
\to B_*>0,
\]
there exists $\mu_0>0$ such that
\[
E(\mu)\geq\frac{B_*}{2}\sqrt{\mu}
\]
for all $0<\mu<\mu_0$.
Evaluating this inequality at
$\mu=\mu_{\mathrm{sc}}(c)$ gives
\[
\mu_{\mathrm{sc}}(c)=O(c^2).
\]
Substituting this estimate into the preceding relation yields
\[
c^2
=
B_*^2\mu_{\mathrm{sc}}(c)
+
O(c^4).
\]
Consequently,
\[
\mu_{\mathrm{sc}}(c)
=
\frac{1}{B_*^2}c^2
+
O(c^4).
\]

Since
\[
\eta_{\mathrm{sc}}(c)
=
\eta_f+\mu_{\mathrm{sc}}(c),
\]
we obtain
\[
\eta_{\mathrm{sc}}(c)-\eta_f
=
K_\diamond c^2+O(c^4),
\]
where
\[
K_\diamond=\frac{1}{B_*^2}.
\]

Using
\[
B_*=\lambda_*A_*\|v_*\|_\diamond
\]
gives
\[
K_\diamond
=
\frac{1}
{\lambda_*^2A_*^2\|v_*\|_\diamond^2}.
\]
Since
\[
A_*^2=\frac{\lambda_*}{\ell_*},
\]
we further obtain
\[
K_\diamond
=
\frac{\ell_*}
{\lambda_*^3\|v_*\|_\diamond^2}.
\]
Finally,
\[
\ell_*=\frac{\Gamma_*}{3\lambda_*},
\]
and therefore
\[
K_\diamond
=
\frac{\Gamma_*}
{3\lambda_*^4\|v_*\|_\diamond^2}.
\]
By {\rm (H4)}, $\Gamma_*>0$, so
\[
K_\diamond>0.
\]
\end{proof}

% ---------------------------------------------------------
\subsection{Explicit tensor form of the separation coefficient}
\label{subsec:explicit-K}
% ---------------------------------------------------------

The coefficient in
Theorem~\ref{thm:quadratic-separation} can be written entirely in
terms of the local differential geometry of the potential.

\begin{corollary}[Explicit separation coefficient]
\label{cor:explicit-K}
Under the assumptions of
Theorem~\ref{thm:quadratic-separation},
\begin{equation}
\boxed{
K_\diamond
=
\frac{
3T_*[v_*,v_*,H_*^{-1}g_*]
-
Q_*[v_*,v_*,v_*,v_*]
}{
3\lambda_*^4
\|v_*\|_\diamond^2
},
}
\label{eq:explicit-K-tensor}
\end{equation}
where
\[
g_*
=
T_*[v_*,v_*,\cdot]^\sharp.
\]

Equivalently, in an orthonormal Hessian eigenbasis
$\{v_i\}_{i=1}^d$,
\begin{align}
K_\diamond
&=
\frac{1}{
3\lambda_*^4
\|v_*\|_\diamond^2
}
\Bigg[
3
\sum_{i=1}^d
\frac{
T_*[v_*,v_*,v_i]^2
}{
\lambda_i
}
\nonumber\\
&\hspace{3em}
-
Q_*[v_*,v_*,v_*,v_*]
\Bigg].
\label{eq:explicit-K-eigenbasis}
\end{align}
\end{corollary}

\begin{proof}
Substitute the representations
\eqref{eq:Gamma-star} and
\eqref{eq:Gamma-eigenbasis}
into \eqref{eq:K-Gamma}.
\end{proof}

\begin{remark}[Geometric interpretation]
\label{rem:K-interpretation}
The coefficient $K_\diamond$ combines three distinct pieces of
local geometry.

First,
\[
\lambda_*
\]
determines the linear flip threshold and the leading displacement of
the gradient in the critical direction.

Second,
\[
\Gamma_*
\]
determines the nonlinear saturation of the unstable critical mode
and hence the amplitude of the emerging period-two orbit.

Third,
\[
\|v_*\|_\diamond
\]
measures the critical eigendirection in the norm used by the
clipping rule.

Thus the switching-contact scale is determined jointly by the
smooth bifurcation geometry and the geometry of the clipping norm.
\end{remark}

% ---------------------------------------------------------
\subsection{Euclidean gradient clipping}
\label{subsec:euclidean-separation}
% ---------------------------------------------------------

For standard gradient clipping,
\[
\|\cdot\|_\diamond=\|\cdot\|_2.
\]
Since
\[
\|v_*\|_2=1,
\]
the separation law simplifies substantially.

\begin{corollary}[Euclidean clipping]
\label{cor:euclidean-separation}
For the standard Euclidean clipping rule,
\begin{equation}
\boxed{
\eta_{\mathrm{sc}}(c)
-
\frac{2}{\lambda_*}
=
K_2c^2
+
O(c^4),
}
\label{eq:euclidean-main-law}
\end{equation}
where
\begin{equation}
\boxed{
K_2
=
\frac{\ell_*}{\lambda_*^3}
=
\frac{\Gamma_*}{3\lambda_*^4}.
}
\label{eq:K-euclidean}
\end{equation}

Equivalently,
\begin{equation}
\boxed{
K_2
=
\frac{1}{3\lambda_*^4}
\left[
3
\sum_{i=1}^d
\frac{
T_*[v_*,v_*,v_i]^2
}{
\lambda_i
}
-
Q_*[v_*,v_*,v_*,v_*]
\right].
}
\label{eq:K-euclidean-expanded}
\end{equation}
\end{corollary}

% ---------------------------------------------------------
\subsection{The same coefficient controls two phenomena}
\label{subsec:coefficient-duality}
% ---------------------------------------------------------

The separation coefficient admits a direct interpretation in terms
of the smooth flip coefficient.

\begin{corollary}[Flip--switching coefficient relation]
\label{cor:flip-switching-relation}
The first flip coefficient $\ell_*$ and the leading
switching-separation coefficient satisfy
\begin{equation}
\boxed{
K_\diamond
=
\frac{\ell_*}{
\lambda_*^3
\|v_*\|_\diamond^2
}.
}
\label{eq:flip-switching-coefficient}
\end{equation}

Therefore, within the present gradient structure, the same nonlinear
coefficient that determines the birth and stability of the smooth
period-two branch also determines its leading-order distance from
the clipping boundary in parameter space.
\end{corollary}

\begin{remark}[Interpretation for Euclidean clipping]
\label{rem:coefficient-duality-euclidean}
For Euclidean clipping,
\[
K_2
=
\frac{\ell_*}{\lambda_*^3}.
\]
Hence a larger positive supercritical flip coefficient produces a
smaller initial oscillation amplitude,
\[
A_*^2=\frac{\lambda_*}{\ell_*},
\]
and therefore delays the first clipping contact to a larger value of
$\eta-\eta_f$.

Conversely, a smaller positive $\ell_*$ produces a larger
small-amplitude cycle for the same $\mu$ and causes clipping to be
encountered sooner.
\end{remark}

% ---------------------------------------------------------
\subsection{Dynamical regimes around the first contact}
\label{subsec:regimes-around-contact}
% ---------------------------------------------------------

The first-contact parameter separates two distinct local regimes.

\begin{proposition}[Pre-contact and contact regimes]
\label{prop:precontact-contact-regimes}
There exist $c_0>0$ and $\mu_0>0$ such that for every
$0<c<c_0$ the following statements hold.

\begin{enumerate}[label=\textnormal{(\roman*)}]

\item
If
\[
0<\mu<\mu_{\mathrm{sc}}(c),
\]
then
\begin{equation}
x_\pm(\mu)\in\Omega_c^-.
\label{eq:precontact-smooth}
\end{equation}
Consequently,
$\mathcal O_2(\mu)$ is simultaneously a stable period-two cycle of
the smooth map $G_{\eta_f+\mu}$ and of the clipped map
$F_{\eta_f+\mu,c}$.

\item
At
\[
\mu=\mu_{\mathrm{sc}}(c),
\]
both points satisfy
\begin{equation}
x_+(\mu_{\mathrm{sc}})
\in\Sigma_c,
\qquad
x_-(\mu_{\mathrm{sc}})
\in\Sigma_c.
\label{eq:both-on-switching}
\end{equation}
The same two points still satisfy the period-two equations for
$F_{\eta_{\mathrm{sc}},c}$.

\item
If
\[
\mu_{\mathrm{sc}}(c)<\mu<\mu_0,
\]
then the continuation of the \emph{smooth reference branch}
satisfies
\begin{equation}
x_\pm(\mu)\in\Omega_c^+.
\label{eq:postcontact-reference-clipped}
\end{equation}
In general, however, this smooth reference cycle is no longer a
period-two cycle of the clipped map.

\end{enumerate}
\end{proposition}

\begin{proof}
By Proposition~\ref{prop:exposure-monotonicity},
the exposure $E(\mu)$ is strictly increasing for all
sufficiently small $\mu>0$. We choose $\mu_0>0$ small enough
that this monotonicity holds throughout $(0,\mu_0)$.

If
\[
0<\mu<\mu_{\mathrm{sc}}(c),
\]
then
\[
E(\mu)<E(\mu_{\mathrm{sc}}(c))=c.
\]
By Theorem~\ref{thm:simultaneous-contact},
both points therefore lie in the smooth region:
\[
x_\pm(\mu)\in\Omega_c^-.
\]
Since $F_{\eta,c}$ and $G_\eta$ coincide at both cycle points,
the smooth period-two cycle is also a period-two cycle of the
clipped map.

At
\[
\mu=\mu_{\mathrm{sc}}(c),
\]
we have
\[
E(\mu_{\mathrm{sc}}(c))=c.
\]
Hence
\[
x_+(\mu_{\mathrm{sc}})\in\Sigma_c,
\qquad
x_-(\mu_{\mathrm{sc}})\in\Sigma_c.
\]
Moreover, the clipping operator satisfies
\[
C_c^\diamond(g)=g
\qquad
\text{whenever }
\|g\|_\diamond=c.
\]
Thus the clipped and smooth maps still agree at both cycle
points, and the period-two relations remain valid at the
contact parameter.

Finally, if
\[
\mu_{\mathrm{sc}}(c)<\mu<\mu_0,
\]
strict monotonicity gives
\[
E(\mu)>c.
\]
Therefore
\[
x_\pm(\mu)\in\Omega_c^+.
\]
These points still form the continuation of the smooth reference
cycle, but in general they no longer form a period-two cycle of
the clipped map, because in $\Omega_c^+$ the gradient update is
rescaled by the factor
\[
\frac{c}{\|\nabla V(x)\|_\diamond}<1.
\]
\end{proof}

\begin{remark}[Stability at the contact point]
\label{rem:contact-stability-not-claimed}
At the exact switching-contact parameter, the two-cycle continues
to exist as an orbit of the clipped map because the map itself is
continuous across $\Sigma_c$.

We do not, however, infer its stability directly from the smooth
multiplier formula. The derivative of the clipped map has different
one-sided limits at the switching surface, so stability at and
beyond contact is a genuinely piecewise-smooth problem.

This issue is deferred to the post-contact analysis.
\end{remark}
\subsection{Origin of the improved remainder}
\label{subsec:origin-improved-remainder}

The $O(c^4)$ remainder in
Theorem~\ref{thm:quadratic-separation}
is a structural consequence of the exact gradient-balance
identity.

Indeed, Theorem~\ref{thm:exposure-expansion} gives
\[
E(\mu)
=
B_*\sqrt{\mu}
+
O(\mu^{3/2}),
\]
rather than the more general expansion
\[
E(\mu)
=
B_1\sqrt{\mu}
+
B_2\mu
+
O(\mu^{3/2}).
\]
The absence of the $O(\mu)$ term is a consequence of the
oddness of the half-difference in the signed square-root
parameter and of the exact identity
\[
E(\mu)
=
\frac{2}{\eta_f+\mu}
\|d(\sqrt{\mu})\|_\diamond .
\]

Squaring the actual exposure expansion gives
\[
E(\mu)^2
=
B_*^2\mu
+
O(\mu^2).
\]
At first contact,
\[
E(\mu_{\mathrm{sc}}(c))=c,
\]
and therefore
\[
c^2
=
B_*^2\mu_{\mathrm{sc}}(c)
+
O\bigl(\mu_{\mathrm{sc}}(c)^2\bigr).
\]
Since
\[
\mu_{\mathrm{sc}}(c)=O(c^2),
\]
it follows that
\[
\mu_{\mathrm{sc}}(c)
=
\frac{1}{B_*^2}c^2
+
O(c^4).
\]

Thus the improvement from the naively expected $O(c^3)$
remainder to $O(c^4)$ does not require a higher-order normal-form
coefficient. It follows directly from the exact two-cycle
gradient balance and the resulting parity structure of the
exposure.

A further expansion of the form
\[
\eta_{\mathrm{sc}}(c)
=
\eta_f
+
K_\diamond c^2
+
K_{\diamond,4}c^4
+
o(c^4)
\]
would require additional higher-order information on the
period-two branch. Such coefficients are not needed for the
leading separation law considered here.